\section{Discussion and Threats to Validity}

\SystemName does not prove current quantum advantage. It proves a computer-architecture point: once task, quality target, and native baseline are fixed, advantage becomes a measured design-space boundary that says whether future effort should buy faster logical gates, more parallel shots, lower error/control overhead, stronger native co-design, or better algorithmic quality.

The largest threat is an underpowered native baseline. We reduce it by selecting the fastest valid method among ML classifiers, exact/heuristic MaxCut solvers, dense/sparse molecular-energy solvers, and dense/sparse Hamiltonian-simulation solvers, then measuring the strong-native shift. The same-input ML production-native gate with PyTorch AMP CNN/MLP and XGBoost reaches native-level quality but does not tighten the median threshold for tiny digits because GPU framework overhead dominates. Stronger natives on larger tasks, including deeper CNN/ResNet families, boosted trees, tuned MILP/domain heuristics, tensor networks, projector QMC, and CCSD(T)-class molecular references, can still move the frontier rightward; the method quantifies how much that stronger native path would raise the quantum requirement. The suite is not exhaustive: molecule cases use small active spaces, HamSim uses small qubit counts, QAOA uses fixed grids, and 4--16 qubit measurements are not deployment-scale extrapolations.

\cuQuantum is instrumentation, not the future hardware stack. Real hardware adds logical gate time, readout, finite shots, feedback, mapping, routing, data loading, queueing, and error correction; hardware-specific studies should replace $T_{\mathrm{error}}$ and $P_{\mathrm{shots}}$ with code distance, cycle time, logical-error budget, magic-state throughput, decoder latency, controller bandwidth, and device concurrency. The FT-stack sensitivity is stylized, but it keeps decoder/control overhead serial and shows where more shot parallelism stops buying advantage. The same record can carry energy by comparing $E_{\mathrm{native}}=N_{\mathrm{GPU}}P_{\mathrm{GPU}}T_{\mathrm{native}}$ with a quantum budget that adds refrigerator, decoder FPGA/ASIC, control, and host power over $T_{\mathrm{qhw}}$; we leave those watts to facility-specific measurements. The scaling plateau is diagnosed from workflow structure, Slurm artifacts, and Nsight/dmon evidence showing 0.8\% GPU-kernel time and 99.2\% host orchestration on a small ML native gate. Nsight Compute counter collection failed because a driver profiling resource was unavailable, so we do not claim Roofline saturation. Bundled Slurm runs, 12 warmup-separated trials, zero failures, maximum quantum-runtime CV 0.0400, and artifact audits reduce timing and provenance risk.
